\subsection{AQ-136: recovered admissibility, a strict counterexample, and the rational comparison}

\paragraph{Original pages and conventions.}
The official cote 136 PDF has alternating mathematical sheets and
unrelated printed reverses.  The relevant physical PDF pages are
111, 113, 115, 117, 119 and 121, with archive footers respectively
110, 112, 114, 116, 118 and 120.  Thus the supplied intake's
``136/117--121'' must not be read as six consecutive mathematical
pages or as the archive-footer numbering.
Physical 111 introduces contravariant functors from simplicial
sets to nonnegatively bigraded commutative differential
$S_k$-algebras with divided powers, and identifies the functors
carrying colimits to limits with simplicial such algebras.
Physical 113--115 requires comparison with constant-coefficient
cohomology after the weight-tail operation $\tau_\infty$.
The coefficient map is $S_k\longrightarrow C_n^{**}$; the sheet
calls this the augmentation.  No map $C_n^{**}\to S_k$ is assumed.

The explicit tests on physical 117 are
\begin{equation}
 \begin{aligned}
 S_k^i&\longrightarrow H^{i,0}(C_n^{**})
  &&\text{is an isomorphism for sufficiently large }i,\\
 H^{i,p}(C_n^{**})&=0
  &&(p>0,\ i\geq\lambda(p)).
 \end{aligned}                                           \tag{136.3}
\end{equation}
The final lines distinguish an initially simplex-dependent
$\lambda_n$ from the desired bound independent of $n$.
On physical 119, let $\pi(i)$ be the smallest $p$ for which the
simplicial module $C_*^{i,p}$ is not flasque, and $+\infty$ if
there is none.  The sheet explains flasqueness by vanishing
simplicial homology, and asks that
\[
                         \pi(i)\longrightarrow+\infty.
                                                               \tag{136.4}
\]
Physical 121 asks for a common admissible source
\[
 C''\longrightarrow C,\qquad C''\longrightarrow C',              \tag{136.a}
\]
and, for parallel morphisms $f,g:C\rightrightarrows C'$, an
admissible $h:C''\to C$ \emph{``qui `\'egalise' ''}, namely
$fh=gh$.  The same sheet also asks whether the groupoid obtained
by inverting all morphisms is connected and simply connected.
The strict equalizer condition is stronger than that groupoid
assertion.  The proofs below exhibit the difference concretely.

\paragraph{Precise category used in the proofs.}
Write $S=\Gamma_k(T)=\bigoplus_{w\geq0}kT^{[w]}$, with
$T^{[a]}T^{[b]}=\binom{a+b}{a}T^{[a+b]}$.
The coefficient weight is $w$, the cochain degree is $p$, and
$d$ has bidegree $(0,1)$.  Let $\mathcal A_k$ consist of simplicial,
nonnegatively bigraded, strictly commutative PD differential
$S$-algebras $C_n^{w,p}$, with morphisms preserving $S$, both
gradings, the differential and divided powers.  Their PD ideal
contains all positive-weight even elements; odd elements square
to zero.  Require (136.3) with bounds uniform in $n$, and
(136.4).  Here flasque means that the normalized simplicial
chain complex is acyclic, including in degree zero.
The associated functor on arbitrary simplicial sets is obtained
by taking compatible sections over all simplices, separately in
each bidegree.  This construction preserves the algebra and PD
operations and is the colimit-to-limit extension on physical 111.
These conventions specify the category rather than presuming
unretrieved prior-run definitions.  The strict counterexample
below needs only the eventual degree-zero case of (136.4), so
it also applies to a larger interpretation with weaker bounds
in (136.3).  The rational assertions concern the explicit
category $\mathcal A_k$ just defined.

\paragraph{Why these explicit tests give the cohomological comparison.}
For completeness, the numerical tests used here do imply the
weight-tail comparison on physical 113--115.  More precisely,
for every fixed $d$ there is $W_d$ such that, for every simplicial
set $X$ and $w\geq W_d$, the two natural comparison maps give
\[
 H^j(C^{w,*}(X))\simeq H^j(X;k)\otimes_k S^w
                  \qquad(0\leq j\leq d).                  \tag{136.2}
\]
This assertion uses the uniform-in-simplex form of (136.3).

\begin{proof}
Let $\mathcal P$ be the abelian category of simplicial
$k$-modules.  Sections over $X$ are
$\Gamma_X(F)=\operatorname{Hom}_{\mathcal P}(k[X],F)$.
The normalized chain complex $K=N(k[X])$ has free $k$-modules,
with basis the nondegenerate simplices.  We first check that
if $F$ has acyclic normalization, then it is acyclic for the
right-derived functors of $\Gamma_X$.

Under Dold--Kan this is the assertion
$\operatorname{Ext}^r_{\operatorname{Ch}_{\geq0}(k)}
 (K,N(F))=0$ for $r>0$.
Choose a resolution of $K$ by projective chain complexes,
using sums of projective disks and degree-zero projectives.
Every short exact sequence in this resolution is split in
each chain degree: the first quotient has projective terms,
and its kernel then has projective terms, inductively.
For any nonnegative chain complex $J$ of projective modules,
the cochain complex $\operatorname{Hom}^{\bullet}(J,N(F))$
is acyclic.  Indeed a closed homogeneous map is
null-homotopic by induction on the source degree: subtract
the already chosen homotopy on the boundary, obtain a map
into cycles, and lift it to the next term of the acyclic
target using projectivity of that source term.  This begins
in degree zero and works for every homogeneous degree.

Now a chain map from a kernel in the projective resolution
to $N(F)$ has a graded extension to the preceding projective
complex.  Its defect of commuting with the differential
factors through the quotient and is a closed degree-one
Hom element.  The just-proved Hom acyclicity corrects this
defect.  Thus every such chain map extends as a chain map;
applying Hom to the resolution is exact in positive degrees.
This proves the required derived-section acyclicity.

The same projective-disk resolution gives
$R^j\Gamma_X(\underline{k})=H^j(X;k)$.
One can verify this without identifying a topological site:
on complexes with projective terms, the cohomology of
$\operatorname{Hom}^{\bullet}(K,k[0])$ has degree-zero
functor $\operatorname{Hom}_{\rm Ch}(K,k[0])$ and vanishes
in positive degrees on projective disks and degree-zero
projectives.  Degreewise split exact sequences give its
long exact sequence, so the projective resolution computes
exactly these groups.  The resulting cochain complex is
the usual normalized simplicial cochain complex of $X$.

Choose $w$ so large that (136.3) holds through cochain
degree $d+1$, and that all rows $C_*^{w,p}$ for
$0\leq p\leq d+1$ are acyclic as simplicial modules.
There is one such bound for every $d$, by (136.3)--(136.4).
The first-quadrant hypercohomology spectral sequence for
$R\Gamma_X(C^{w,*})$, applied first to vertical cohomology,
identifies its groups through degree $d$ with
$R^j\Gamma_X(\underline{S^w})$.
Applied instead to the derived sections of each term,
it identifies those same groups with
$H^j(\Gamma_X(C^{w,*}))$, by the acyclicity just proved.
The coefficient augmentation and the map from ordinary
to derived sections are the comparison maps in these two
identifications.  Since $S^w$ is a rank-one free
$k$-module, they give (136.2).
Both spectral sequences converge in these degrees:
all indices are nonnegative and each total degree uses
only finitely many bidegrees.  This proves the claim for
arbitrary $X$, without interchanging a weight-tail limit
with an infinite product.
\end{proof}

\paragraph{Simplicial algebra fact used below.}
Let $L_n=k^{n+1}$ have basis $x_0,\ldots,x_n$.  For
$\theta:[m]\to[n]$, put
\[
 \theta^*x_j=\sum_{\theta(a)=j}x_a.
\]
Then $t_n=\sum_jx_j$ defines an inclusion of the constant
simplicial module $k$ in $L$, and put $Q=L/kt$.
The normalized chain complex of $L$ is the two-term disk
$k\xrightarrow{\pm1}k$ in degrees one and zero: in the
normalization using $\bigcap_{i<n}\ker d_i$, its degree-one
term is $kx_0$, and in degrees at least two the simultaneous
kernels of $d_0$ and $d_1$ are zero.  Thus $L$ is simplicially
homotopy equivalent to zero.  An exact known input for the last
step is the chain-to-simplicial homotopy theorem, Stacks Project,
Lemma 14.29.3, Tag \texttt{01A4}.

For every $r>0$, $\Gamma^r(L)$ and $\operatorname{Sym}^r(L)$
are consequently simplicially contractible: apply the functor
to a homotopy between the identity and zero.  Functorial
preservation of simplicial homotopies is Stacks Project,
Remark 14.26.4, Tag \texttt{019M}; these positive-degree
functors send the zero map to zero.  If $V$ is a simplicially
contractible module and $W$ any simplicial module, then
$V\otimes_k W$ is also contractible.  Indeed the homotopy
maps $h_{n,\alpha}:V_n\to V_n$, indexed by
$\alpha:[n]\to[1]$, give homotopy maps
$h_{n,\alpha}\otimes1$ on $V_n\otimes W_n$, satisfying the
same face and degeneracy identities.  No flatness or
characteristic-zero hypothesis is used in this observation.

\paragraph{AQ-136-DP-MODEL: an explicit admissible integral theory.}
For every commutative ring $k$, define
\[
 B_n=\Gamma_k(x_0,\ldots,x_n)
       \otimes_k\Lambda_k(e_0,\ldots,e_n)/(e_0+\cdots+e_n),
 \quad dx_j=e_j,\quad de_j=0,\quad T\mapsto\sum_jx_j .        \tag{136.5}
\]
Each $x_j$ has bidegree $(1,0)$ and $e_j$ has bidegree
$(1,1)$.  Simplicial operators use the displayed summation
formula on both $x_j$ and $e_j$.  Divided powers of the
coefficient $T$ are sent to divided powers of its sum.
Then $B\in\mathcal A_k$; (136.3) holds in every weight,
and its row $B^{w,p}_*$ is contractible whenever $w>p$.

\begin{proof}
First, all the displayed algebraic structures exist integrally.
The ordinary divided-power polynomial algebra has basis
$\prod_jx_j^{[a_j]}$ and the multiplication above; its precise
universal property is Stacks Project, Lemma 23.5.1,
Tag \texttt{07H5}.  Tensoring it with an exterior algebra
gives the usual strictly graded-commutative PD algebra.
One direct verification is to work first over $\mathbb Z$.
It is a free abelian group inside its rationalization, where
$\gamma_r(z)=z^r/r!$ for positive-weight even $z$.
For an even monomial containing an exterior generator, its
square is zero, so its divided powers of order at least two
vanish.  For a positive polynomial divided-power monomial,
$r!$ divides the factorial coefficient in its $r$th power:
choose one variable occurring positively and count partitions
of its $ra$ occurrences into $r$ unordered blocks of size $a$.
The addition formula for divided powers therefore preserves
the integral lattice.  The identities hold in the
rationalization and hence over $\mathbb Z$; the resulting
integer polynomial identities give the construction over every
$k$, including rings with torsion.  Eliminating $e_0$ gives
an exterior algebra on $e_1,\ldots,e_n$, so the displayed
relation is compatible with this structure.

The differential is explicitly
$d(x_j^{[a]})=x_j^{[a-1]}e_j$ for $a>0$ and extends by the
signed Leibniz rule.  Its square is zero.  It satisfies
$d\gamma_r(z)=\gamma_{r-1}(z)dz$: this follows over
$\mathbb Q$ by differentiation and integrally by the same
free-lattice argument.  The simplicial maps preserve it.
The relation $\sum e_j=0$ says precisely that $dT=0$.
All simplicial identities follow from composition of the
sums over the fibres of order-preserving maps.

At a fixed simplex make the invertible linear change of
variables $(x_0,\ldots,x_n)\mapsto(T,x_1,\ldots,x_n)$.
The PD universal property gives
\[
 B_n\simeq S\otimes_k\bigotimes_{j=1}^n
            \bigl(\Gamma_k(x_j)\otimes\Lambda_k(e_j)\bigr).
                                                               \tag{136.6}
\]
For one factor a contraction onto $k$ is
$h(x^{[a]}e)=x^{[a+1]}$ and $h(x^{[a]})=0$.
The identity $dh+hd=1-\epsilon$ follows on each basis
monomial; $\epsilon$ keeps only the constant term.
The tensor contraction is obtained successively by
$h\otimes1+\epsilon\otimes h$, with the usual sign in the
tensor action of a degree-minus-one map.  Thus the inclusion
$S\to B_n$ is a cochain homotopy equivalence.  The contractions
preserve weight, so $H^{w,0}(B_n)=S^w$ and all
$H^{w,p}(B_n)$ for $p>0$ vanish, uniformly in $n$.

Finally there is a natural identification of simplicial modules
\[
 B_*^{w,p}=\Gamma^{w-p}(L)\otimes_k\Lambda^p(Q),             \tag{136.7}
\]
with zero understood for $w<p$.  If $w>p$, the first factor
is contractible by the preceding simplicial fact, and so is
the tensor product.  Hence (136.4) holds with
$\pi(w)\geq w$.  This proves every assertion.
\end{proof}

\paragraph{AQ-136-STRICT: the original strict equalizing-source conjecture is false.}
If $k\ne0$, there are two parallel morphisms in
$\mathcal A_k$ which admit no equalizing morphism from any
admissible source.  In particular this refutes the strict
question on physical 121 over $\mathbb Z$, without a
rationalization or a proposed lifting lemma.
The supplied intake had reported a strict-form refutation
under the earlier identifier C07; the following is a complete
independent proof from the recovered original conditions,
not a claim to have recovered that absent proof.

\begin{proof}
Use the theory $B$ of (136.5).  Define a differential
$B$-module $M$ by adjoining symbols $\alpha_j$ of bidegree
$(1,0)$ and $\beta_j$ of bidegree $(1,1)$, imposing
$\sum\alpha_j=\sum\beta_j=0$, and putting
$d\alpha_j=\beta_j$, $d\beta_j=0$.
In invariant form
\[
 M=B\otimes_k\bigl(Q\{0\}\xrightarrow{1}Q\{1\}\bigr),
                                                               \tag{136.8}
\]
where the two copies have coefficient weight one and the
braces indicate their cochain degrees.  There is a
weight-preserving degree-zero PD derivation
\[
 \delta:B\longrightarrow M,\qquad
 \delta x_j=\alpha_j,\quad\delta e_j=\beta_j,\quad
 \delta(x_j^{[a]})=x_j^{[a-1]}\alpha_j .                    \tag{136.9}
\]
Extend it by the degree-zero Leibniz rule.  The two sum
relations make $\delta T=0$ and $\delta\sum e_j=0$.
The divided-power addition formula shows that it annihilates
all of $S$ and respects the PD identities.  It commutes
with $d$, as is checked on $x_j,e_j$ and their divided
powers and then by the Leibniz rule.  Everything is natural
in the simplex.

Form the square-zero extension $D=B\ltimes M$, with
$M^2=0$, its induced differential, and divided powers
\[
 \gamma_r(b+m)=\gamma_r(b)+\gamma_{r-1}(b)m
 \quad(r\geq1)                                             \tag{136.10}
\]
for even elements in the PD ideal.  Give $M$ zero higher
divided powers.  Formula (136.10), the addition formula,
and $M^2=0$ verify all the PD axioms; alternatively these
are the terms of degree zero and one in an extra
square-zero variable in the integral construction just
proved.  Differential compatibility follows because
\[
 d(\gamma_{r-1}(b)m)
  =\gamma_{r-2}(b)(db)m+\gamma_{r-1}(b)dm,
\]
with the first term absent for $r=1$, agreeing with
$\gamma_{r-1}(b+m)d(b+m)$ after using $M^2=0$.
Thus
\[
 j_0,j_1:B\rightrightarrows D,
 \qquad j_0(b)=(b,0),\quad j_1(b)=(b,\delta b)              \tag{136.11}
\]
are strict morphisms of differential PD $S$-algebras.

We verify that $D$ is admissible.  The summand $M$ is
vertically contractible: on homogeneous $b\in B$ put
$h(b\beta_j)=(-1)^{\deg b}b\alpha_j$ and
$h(b\alpha_j)=0$.  The relations are preserved, and
$dh+hd=1$ follows by cancelling the two terms involving
$db$.  Consequently $S\to D_n$ has exactly the same
cohomology as $S\to B_n$ in every weight and degree.
Horizontally,
\[
 M_*^{w,p}=
  (B_*^{w-1,p}\otimes Q)\ \oplus\
  (B_*^{w-1,p-1}\otimes Q).                                \tag{136.12}
\]
The second term is zero if $p=0$.  Both terms are
contractible for $w>p+1$, by (136.7) and the tensor
homotopy fact.  Hence $\pi_D(w)\geq\max(w-1,0)$.
No coefficient tail has been discarded: both $B$ and
$D$ retain every $T^{[w]}$ over the original ring $k$.

At every simplex the equalizer of $j_0,j_1$ is exactly
$S$.  To check this, eliminate the zeroth variable as
in (136.6); $M_n$ is free over $B_n$ on
$\alpha_1,\ldots,\alpha_n,\beta_1,\ldots,\beta_n$.
Expand an element of $B_n$ uniquely as
\[
 \sum_{a,I}s_{a,I}x_1^{[a_1]}\cdots x_n^{[a_n]}e_I,
 \qquad s_{a,I}\in S.
\]
The coefficient of $\beta_j$ in its $\delta$ is the
signed exterior partial derivative with respect to $e_j$.
The basis monomials containing $e_j$ map injectively to
distinct monomials after removing it.  Thus $\delta=0$
forces every nonempty $I$ coefficient to vanish.
Then the coefficient of $\alpha_j$ is the divided-power
partial derivative, sending $x_j^{[a_j]}$ to
$x_j^{[a_j-1]}$ with coefficient one.  It forces all
$a_j>0$ coefficients to vanish.  This argument works
over arbitrary $k$; it does not divide by an integer.
The remaining elements are exactly $S$.

Suppose an admissible $E$ and an $S$-morphism $u:E\to B$
equalized (136.11).  By the computed equalizer, $u$
factors through the constant simplicial algebra $S$.
Since it is an $S$-morphism, this factorization retracts
the structure map $S\to E$.  In each weight $w$ it makes
the nonzero constant simplicial module
$kT^{[w]}$ a direct summand of $E_*^{w,0}$.
Therefore $H_0^{\rm simp}(E_*^{w,0})$ contains $k$ for
every $w$.  This contradicts (136.4), which makes that
row acyclic for all sufficiently large $w$.
The contradiction proves the strict assertion false.
\end{proof}

\paragraph{AQ-136-RATIONAL: common sources and terminal localization over a $\mathbb Q$-algebra.}
If $k$ is a $\mathbb Q$-algebra, every finite family of
objects of $\mathcal A_k$ has a common source in
$\mathcal A_k$.  The ordinary category obtained from
$\mathcal A_k$ by inverting every morphism is equivalent
to the terminal groupoid.  This is a statement about
that groupoid, not about all higher homotopy groups of
an infinity-categorical localization.

\begin{proof}
Identify $S=\Gamma_k(T)$ with $k[t]$ by
$T^{[w]}=t^w/w!$.  For an integer $N\geq1$, define
\[
 B_{N,n}=
 \frac{k[t,x_0,\ldots,x_n]\otimes\Lambda(e_0,\ldots,e_n)}
      {(\sum x_j-t^N,\ \sum e_j)},
 \quad dx_j=e_j,\quad dt=de_j=0.                           \tag{136.13}
\]
Put $\operatorname{wt}(t)=1$ and
$\operatorname{wt}(x_j)=\operatorname{wt}(e_j)=N$.
Use the same simplex operators as before, fixing $t$.
Over a $\mathbb Q$-algebra the PD operations on even
positive-weight elements are uniquely $z^r/r!$.
Thus (136.13) is a differential PD $S$-algebra.

Eliminating $x_0,e_0$ gives the polynomial de Rham
complex over $S$ in $n$ variables of weight $N$.
In one variable the homotopy
$h(x^ae)=x^{a+1}/(a+1)$ contracts onto the constants.
Its tensor product contraction proves (136.3) in every
weight, again preserving weight.  For the horizontal
calculation let $A$ be the same construction with
coefficient $U=\sum x_j$ and all $x_j,e_j,U$ of weight
one.  As a simplicial module
$A_*^{m,p}=\operatorname{Sym}^{m-p}(L)\otimes\Lambda^p(Q)$,
so it is contractible if $m>p$.
Before extension of coefficients regrade these weights
by $N$.  We have
\[
 B_N=k[t]\otimes_{k[U]}A,\qquad U\longmapsto t^N,
 \qquad k[t]=\bigoplus_{r=0}^{N-1}t^r k[U].                \tag{136.14}
\]
The last equality is an actual free-module decomposition,
valid over every $k$.  The component of weight $w$
therefore identifies with $A^{m,p}$, where
$w=Nm+r$ and $0\leq r<N$.  It is contractible for
$m>p$.  In particular
$\pi_{B_N}(w)\geq\lfloor w/N\rfloor$, which tends to
infinity.  This proves $B_N\in\mathcal A_k$.

Now fix finitely many admissible theories $C_a$.
Choose $N\geq1$ large enough that every simplicial
module $C_a^{N,0}$ is acyclic, using (136.4).
The coefficient $t^N$ determines a map from the
constant simplicial module $k$ to $C_a^{N,0}$.
It extends along the inclusion $k\to L$ to a
simplicial linear map
\[
                         L\longrightarrow C_a^{N,0}.       \tag{136.15}
\]
Here is the precise lifting reason.  Normalization
sends $k\to L$ to the inclusion of the degree-zero
term into the two-term disk.  Vanishing of $H_0$
says that the coefficient element $t^N$ at simplex
zero is the boundary of an element in normalized
degree one.  Choosing that element gives a chain
map from the disk.  The Dold--Kan equivalence,
Stacks Project Theorem 14.24.3, Tag \texttt{019G},
lifts it to (136.15), retaining the given constant
map.  No lift of an arbitrary cochain complex
through a surjection is being asserted.

Send $x_j$ to the image of its basis vector in
(136.15), and $e_j$ to the vertical differential
of that image.  Their sums are respectively
$t^N$ and zero.  The polynomial/exterior universal
property defines a simplicial differential algebra
map $B_N\to C_a$, and it preserves divided powers
because $k$ contains $\mathbb Q$.  It preserves
weights and $S$ by construction.  These maps give
the claimed simultaneous common source.

It remains to prove terminality after inversion;
strict equalizers cannot be used, by the preceding
counterexample.  For any admissible $D$, let
\[
 P(D)_n=D_n\otimes_k k[s,ds],\quad
 \operatorname{wt}(s)=\operatorname{wt}(ds)=0,
 \quad d s=ds,\quad d(ds)=0.                              \tag{136.16}
\]
It has cochain-degree-zero $s$ and degree-one $ds$.
This is an object of $\mathcal A_k$.  Integration of
polynomials contracts $k[s,ds]$ onto $k$, naturally
in $D$, and so leaves every vertical cohomology
unchanged.  Its row of bidegree $(w,p)$ is the sum
of copies of $D^{w,p}$ and $D^{w,p-1}$ indexed by
powers of $s$.  Normalization and homology commute
with these direct sums; hence (136.4) is retained.
All PD operations exist canonically over $\mathbb Q$.
The maps $j:D\to P(D)$ and
$e_0,e_1:P(D)\to D$ are admissible morphisms and
$e_0j=e_1j=1$.  Consequently $e_0=e_1=j^{-1}$ in
the all-arrow localization.

For any parallel $f,g:C\rightrightarrows D$, choose
a common-source map $a:B_N\to C$ as above.
There is an explicit PD differential map
$H:B_N\to P(D)$ given on each generator by
\[
 H(x_j)=(1-s)fa(x_j)+sga(x_j),\qquad H(e_j)=dH(x_j),
 \qquad H(t)=t.                                           \tag{136.17}
\]
Its sums are $t^N$ and zero, it is natural in the
simplex and preserves weights.  Thus the relations
and all structures are respected.  Its two
evaluations are $fa$ and $ga$.  The equality of
evaluations in the localization, followed by
inversion of $a$, gives $f=g$ there.

Finally any finite zigzag in the localization
involves finitely many objects $C_0,\ldots,C_r$.
Choose a common source $B_M$ and arrows
$a_i:B_M\to C_i$.  For every original arrow
$f:C_i\to C_j$ in the zigzag, the two parallel
arrows $fa_i$ and $a_j$ have the same image in
the localization, by the preceding paragraph.
Hence the image of $f$ is $a_j a_i^{-1}$, and
its inverse is $a_i a_j^{-1}$.  The whole zigzag
telescopes to $a_r a_0^{-1}$.  Applying the same
argument to the union of any two finite zigzags
shows that they agree.  Existence of common sources
makes the groupoid connected; $B_1$ makes it
nonempty.  It therefore has exactly one morphism
between each two objects, proving the assertion.
\end{proof}

\paragraph{AQ-136-INTEGRAL-TAIL: the direct one-generator integral analogue fails.}
The strict counterexample retained every coefficient
$T^{[w]}$ and works over $\mathbb Z$.  There is also a
precise obstruction to extending the rational common-source
construction by replacing $t^N$ with $T^{[N]}$.
For $N\geq2$, form the analogous divided-power theory
$B_N^{\rm PD}$ over $S=\Gamma_{\mathbb Z}(T)$ by
\[
 B_{N,n}^{\rm PD}
  =S\langle x_0,\ldots,x_n\rangle
       \otimes\Lambda(e_0,\ldots,e_n)
          /(\sum x_j-T^{[N]},\ \sum e_j),                   \tag{136.18}
\]
where the first relation means the PD quotient, each
$x_j,e_j$ has weight $N$, and $dx_j=e_j$.
Eliminating the zeroth variables gives a free PD
polynomial de Rham complex over $S$, so vertical
acyclicity is retained.  Nevertheless (136.4) fails:
its degree-zero simplicial homology is nonzero in
arbitrarily large weights.

\begin{proof}
In simplex zero the zero-form algebra is $S$; in
simplex one it is $S\langle x_1\rangle$, with the two
faces evaluating $x_1$ at $T^{[N]}$ and at zero.
Therefore its simplicial degree-zero homology is
\[
 H_0^{\rm simp}((B_N^{\rm PD})^{*,0})=S/J_N,
 \qquad J_N=(\gamma_m(T^{[N]}):m\geq1).                    \tag{136.19}
\]
This is a direct calculation of the cokernel of
$d_0-d_1$, using the $S$-basis $x_1^{[m]}$; it does
not presume flat base change of a simplicial
resolution.  The ideal $J_N$ is the full PD ideal
generated by the chosen coefficient, not merely
its ordinary principal ideal: the iterated-PD and
addition identities show that the displayed
ordinary ideal is PD-stable and is the smallest
such ideal containing $T^{[N]}$.

Let $p>N$ be prime.  Only $mN\leq p$ contributes
in weight $p$.  The corresponding product is
\[
 T^{[p-mN]}\gamma_m(T^{[N]})
   =\frac{p!}{m!(N!)^m(p-mN)!}\,T^{[p]}.                   \tag{136.20}
\]
The coefficient is an integer by the PD identities.
Its numerator has exactly one factor $p$, and
none of the factorials in its denominator has
one: $N<p$, $m\leq p/N<p$, and $p-mN<p$.
Thus it is divisible by $p$.
Every weight-$p$ coefficient of $J_N$ is consequently
a multiple of $p$, so $(S/J_N)^p$ surjects onto
$\mathbb Z/p$ and is nonzero.  Primes $p>N$ are
arbitrarily large, proving the failure of (136.4).
For $N=1$ the PD ideal is the entire positive
coefficient ideal; the counterargument does not
apply.  That value would require a target's
weight-one row to be acyclic, which general
admissibility does not guarantee.
\end{proof}

The rational theorem thus does not extend by the
single-coefficient substitution (136.18), even
though its vertical Poincare lemma is valid.
This is a counterexample to that particular
construction, not to the integral common-source
or integral terminal-localization assertions.
Those two integral assertions remain distinct
open gates in this module.  Their original
defining sheets are now recovered, so the gate
is no longer the absence of the original question.
